\section{Introdução}

Este documento apresenta os requisitos para o desenvolvimento de um
sistema de análise e seleção de carteiras de investimentos baseado no
modelo de seleção de carteiras de Markowitz.

A proposta do sistema surgiu da necessidade de auxiliar o usuário na
análise de diferentes combinações de ativos financeiros, considerando
a relação entre risco e retorno. A partir dos dados dos ativos
selecionados, o sistema deverá permitir a avaliação de diferentes
composições de carteira e a identificação de alternativas eficientes
segundo o modelo adotado.

Para definir o sistema de forma clara, o documento parte das
necessidades do usuário e avança para a definição do sistema e para a
especificação dos requisitos que deverão ser atendidos pelo software.
Essa organização permite separar as necessidades e objetivos do
sistema das funcionalidades e restrições que deverão ser consideradas
durante seu desenvolvimento.

O documento está organizado em quatro blocos principais:
\begin{itemize}
\item \textbf{Requisitos do Usuário:} apresentam as necessidades
do usuário e os serviços que ele espera obter com o sistema;


\item \textbf{Definição do Sistema:} apresenta o objetivo, o
escopo, o perfil do usuário, os principais conceitos do domínio,
as regras gerais e as limitações do sistema;

\item \textbf{Especificação de Requisitos de Software:} detalha
as funcionalidades que o software deverá oferecer e as
condições de operação que deverá atender;

\item \textbf{Projeto do Software:} apresenta como os requisitos
serão traduzidos em uma estrutura de software, considerando os
conceitos do domínio, o modelo de Markowitz e a organização dos
componentes do sistema.

\end{itemize}
